\documentclass{article}

% NeurIPS 2024 style
\usepackage[final]{neurips_2024}

\usepackage[utf8]{inputenc}
\usepackage[T1]{fontenc}
\usepackage{hyperref}
\usepackage{url}
\usepackage{booktabs}
\usepackage{amsfonts}
\usepackage{amsmath}
\usepackage{amssymb}
\usepackage{nicefrac}
\usepackage{microtype}
\usepackage{xcolor}
\usepackage{graphicx}
\usepackage{algorithm}
\usepackage{algpseudocode}
\usepackage{bm}
\usepackage{mathtools}
\usepackage{subfigure}

% Convenience macros
\newcommand{\R}{\mathbb{R}}
\newcommand{\E}{\mathbb{E}}
\newcommand{\dd}{\,\mathrm{d}}
\newcommand{\dW}{\,\mathrm{d}W}
\newcommand{\kB}{k_{\mathrm{B}}}
\newcommand{\EES}{\textsc{EES25}}
\newcommand{\RevAdj}{\textsc{ReversibleAdjoint}}
\newcommand{\diffrax}{\texttt{diffrax}}
\newcommand{\EANN}{\textsc{EANN}}
\newcommand{\MBAR}{\textsc{MBAR}}

\title{%
  Differentiable Molecular Dynamics with Explicit Effectively Symmetric Integrators:\\
  Fitting Infrared Spectra at Constant Memory Cost%
}

\author{%
  Luke Thompson\\
  School of Chemistry, University of Sydney\\
  \texttt{luke.thompson@sydney.edu.au}
}

\begin{document}

\maketitle

\begin{abstract}
Differentiable molecular dynamics (MD) enables the end-to-end optimisation of
force-field parameters by back-propagating through entire trajectories.  The
standard approach stores all intermediate states during the forward pass,
requiring $\mathcal{O}(N)$ memory in the trajectory length---a severe bottleneck
for the long simulations needed to resolve low-frequency spectral features.
We benchmark \EES{} (Explicit Effectively Symmetric Runge--Kutta, order 5),
combined with the \RevAdj{} algorithm from \diffrax, as a drop-in replacement
for conventional integrators in an infrared (IR) spectrum fitting task.
\EES{} is a time-reversible explicit Runge--Kutta method; reversibility
enables the adjoint to \emph{reconstruct} trajectory states on the fly during
the backward pass, reducing adjoint memory to $\mathcal{O}(1)$ at a
computational overhead of roughly two forward passes.
We apply this scheme to the problem of fitting an embedded-atom neural network
(\EANN) water force field to experimental liquid-phase IR spectra via
differentiable NVE trajectories, with an \MBAR{}-reweighted density
constraint.  A Langevin NVT variant additionally recasts the stochastic
differential equation as a deterministic rough-path ODE so that
\RevAdj{} can be applied without modification.
Benchmarks on a 192-atom liquid water system demonstrate that
\EES{}+\RevAdj{} matches gradient quality of direct adjoint baselines
(Heun, Tsit5) while using constant XLA scratch memory irrespective of
trajectory length, confirming its suitability for large-scale spectral
fitting.
\end{abstract}

% ─────────────────────────────────────────────────────────────────────────────
\section{Introduction}
\label{sec:intro}

Matching an empirical or machine-learned inter-atomic potential to dynamical
observables---rather than static energies and forces---is essential for
capturing nuclear quantum effects, anharmonicity, and many-body correlations
that shape condensed-phase spectra.  The infrared absorption lineshape of
liquid water, in particular, spans three distinct frequency regions (librational
$<$\,1000\,cm$^{-1}$, bending $\sim$1640\,cm$^{-1}$, and OH-stretch
$\sim$3400\,cm$^{-1}$) that place stringent, complementary demands on the
potential energy surface (PES).

Differentiable programming allows gradients of a spectral loss to flow back
through the entire trajectory, simultaneously adjusting all force-field
parameters.  The key computational challenge is the memory cost of the adjoint
method: storing every intermediate state for a trajectory of $N$ steps demands
$\mathcal{O}(N)$ memory, which quickly becomes prohibitive for trajectories
long enough to converge IR autocorrelation functions.

\paragraph{Contributions.}
\begin{itemize}
  \item We adapt the \EES{} low-storage Runge--Kutta solver to differentiable
    molecular dynamics, enabling $\mathcal{O}(1)$-memory gradient computation
    via \RevAdj{}.
  \item We introduce a \emph{rough-path} re-casting of the Langevin SDE as a
    deterministic, time-reversible ODE, making \RevAdj{} directly applicable
    to NVT (constant temperature) trajectories without approximation.
  \item We provide a modular benchmark comparing \EES{}+\RevAdj{} against
    Heun+\textsc{DirectAdjoint} and Tsit5+\textsc{DirectAdjoint} in terms of
    XLA scratch memory, wall-clock time, and training-loss behaviour.
  \item We integrate the scheme into a complete IR-spectrum fitting pipeline
    that combines differentiable NVE trajectories, a DeepDipole model,
    Filon-FFT autocorrelation, and \MBAR{}-reweighted density regularisation.
\end{itemize}

% ─────────────────────────────────────────────────────────────────────────────
\section{Background}
\label{sec:background}

\subsection{Differentiable Trajectory Optimisation}

We seek force-field parameters $\theta$ that minimise a loss
$\mathcal{L}(\theta)$ defined on a molecular trajectory
$\bm{x}_{0:T} = \{\bm{r}(t), \bm{v}(t)\}_{t=0}^{T}$.
Training requires the gradient $\nabla_\theta \mathcal{L}$, which by the chain
rule is
\begin{equation}
  \nabla_\theta \mathcal{L}
  = \int_0^T
      \underbrace{\frac{\partial \mathcal{L}}{\partial \bm{x}(t)}}_{\text{loss sensitivity}}
      \cdot
      \underbrace{\frac{\partial \bm{x}(t)}{\partial \theta}}_{\text{param sensitivity}}
    \dd t.
\end{equation}
The sensitivity $\partial \bm{x}/\partial\theta$ satisfies the \emph{adjoint}
ODE, whose numerical integration in reverse time requires access to the forward
trajectory $\bm{x}(t)$.

\paragraph{Direct adjoint (checkpointing).}
The simplest approach stores \emph{all} $N$ forward states and replays them
during the backward pass: memory $\mathcal{O}(N)$, computation $\mathcal{O}(N)$.

\paragraph{Reversible adjoint.}
If the forward map $\Phi_\theta : \bm{x}_n \mapsto \bm{x}_{n+1}$ is
time-reversible (i.e.\ $\Phi_\theta^{-1}$ can be computed from $\bm{x}_{n+1}$
alone), then each forward state can be reconstructed on-the-fly during the
backward pass: memory $\mathcal{O}(1)$, computation $\mathcal{O}(N)$ with a
constant factor of roughly 2.

\subsection{Low-Storage Reversible Runge--Kutta Methods}
\label{sec:ees_background}

An explicit Runge--Kutta method is \emph{effectively symmetric} if its
Butcher tableau satisfies
\begin{equation}
  b_i = b_{s+1-i}, \quad
  a_{s+1-i,\,s+1-j} = b_j - a_{ji},
  \qquad i,j = 1,\ldots,s,
\end{equation}
which implies the numerical flow is palindromic and hence time-reversible
to the order of the method.  \EES{} is a two-register, 5th-order method of
this class from the \texttt{diffrax\_lowstorage} library
\cite{diffrax_lowstorage}; it requires only two auxiliary registers beyond
the state itself, making its storage footprint comparable to a fixed-step
Euler method.

% ─────────────────────────────────────────────────────────────────────────────
\section{Physical System}
\label{sec:system}

\subsection{Liquid Water Model}

We simulate a periodic box of $N_\mathrm{mol} = 64$ water molecules
($N_\mathrm{atom} = 192$ atoms) at ambient conditions: $T = 298.15$\,K,
$P = 1.0$\,atm.  Atomic masses are $m_\mathrm{O} = 15.999\,\mathrm{g/mol}$
and $m_\mathrm{H} = 1.008\,\mathrm{g/mol}$.  All lengths are in nanometres
and times in picoseconds (GROMACS unit convention).

\subsection{EANN Interatomic Potential}
\label{sec:eann}

The potential energy $U(\bm{r};\theta)$ is provided by an Embedded-Atom Neural
Network (\EANN) \cite{eann2020} implemented in JAX.  For each atom $i$ of
element $\alpha(i)$, a feature vector $\bm{g}_i \in \R^{d}$ is constructed
from pairwise Gaussian-type orbital (GTO) basis functions.

\paragraph{Radial basis.}
For a pair $(i,j)$ with separation $r_{ij}$, the GTO amplitudes are
\begin{equation}
  \phi_k^{\alpha}(r_{ij})
    = \exp\!\bigl(-\zeta_k^\alpha (r_{ij} - \mu_k^\alpha)^2\bigr)\,
      f_c(r_{ij}),
  \qquad k = 1,\ldots,n_\mathrm{GTO},
\end{equation}
where $\mu_k^\alpha$ and $\zeta_k^\alpha$ are learnable shifts and exponents
per element $\alpha$, and $f_c(r) = \bigl(\tfrac{1}{2}\cos(\pi r/r_c) +
\tfrac{1}{2}\bigr)^2$ is a smooth cutoff at $r_c = 0.6$\,nm.  The
hyperparameters are $n_\mathrm{GTO} = 12$, initial spacing
$\Delta\mu = r_c/(n_\mathrm{GTO}-\tfrac{2}{3}) \approx 0.055$\,nm, and
initial exponent $\zeta = 0.2/\Delta\mu^2 \approx 66$\,nm$^{-2}$.

\paragraph{Angular channels.}
Features are accumulated over angular momentum channels $L = 0, 1, 2$
(spherical, vector, and symmetric-traceless-tensor components of the
displacement unit vector $\hat{\bm{r}}_{ij}$):
\begin{equation}
  \bm{g}_i^{(L)}
    = \left\|
        \sum_{j \neq i}
          w^\alpha_j\,
          Y^{(L)}(\hat{\bm{r}}_{ij})\,
          \bm{\phi}^{\alpha(j)}(r_{ij})
      \right\|^2_F,
  \qquad
  \bm{g}_i = \bigl[\bm{g}_i^{(0)},\,\bm{g}_i^{(1)},\,\bm{g}_i^{(2)}\bigr]
  \in \R^{3 n_\mathrm{GTO}},
\end{equation}
where $\|\cdot\|_F$ denotes the Frobenius (entry-wise $\ell_2$) norm and
$w^\alpha_j$ is an element-dependent weight for atom $j$.

\paragraph{Network.}
Each atom's energy is computed by a two-hidden-layer network with 64 neurons
per layer, gated SiLU activations, and per-layer LayerNorm:
\begin{align}
  \bm{h}^{(0)} &= \bm{W}^{(0)}_\alpha \bm{g}_i + \bm{b}^{(0)}_\alpha, \\
  \bm{h}^{(\ell+1)} &= \mathrm{LN}\!\bigl(\bm{h}^{(\ell)}\bigr)
    \odot \sigma_g\!\bigl(\mathrm{LN}(\bm{h}^{(\ell)})\bigr),
  \quad \ell=0,1, \\
  \epsilon_i &= \bm{w}_\alpha^\top \bm{h}^{(2)} + \epsilon_0,
\end{align}
where $\mathrm{LN}$ is layer normalisation and
$\sigma_g(\bm{h}) = \bm{h} \odot \mathrm{SiLU}(\bm{h})$ is the gated
activation.  The total energy is $U = \sum_i \epsilon_i$.

% ─────────────────────────────────────────────────────────────────────────────
\section{Dynamics and Integration}
\label{sec:dynamics}

\subsection{NVE Trajectories for Spectral Observables}
\label{sec:nve}

For computing the IR spectrum, we run constant-energy (NVE) trajectories
initialised from an NVT-equilibrated ensemble.  The equations of motion are
Newton's second law
\begin{equation}
  \dot{\bm{r}} = \bm{v}, \qquad
  m_i \dot{v}_i = F_i(\bm{r};\theta) \coloneqq -\nabla_{r_i} U(\bm{r};\theta).
\end{equation}

\paragraph{Leapfrog Verlet (baseline).}
The standard leapfrog scheme is a time-reversible, symplectic 2nd-order
integrator:
\begin{align}
  \bm{F}_n &= -\nabla U(\bm{r}_n;\theta), \\
  \bm{v}_{n+1} &= \bm{v}_n + \frac{\bm{F}_n}{m}\,\Delta t, \\
  \bm{r}_{n+1} &= \bm{r}_n + \bm{v}_{n+1}\,\Delta t.
\end{align}
Centre-of-mass momentum is projected out at each step, and periodic boundary
conditions (PBC) are enforced by wrapping molecule centroids into the simulation
box while preserving covalent geometry.

\paragraph{EES25 (proposed).}
We replace the leapfrog kernel with \EES{} acting on the augmented state
$\bm{y} = (\bm{r}, \bm{v})$, treating Newtonian mechanics as the ODE
$\dot{\bm{y}} = f(\bm{y};\theta) = (\bm{v},\, \bm{F}/m)$:
\begin{equation}
  \bm{y}_{n+1} = \EES{}(\bm{y}_n;\,\theta,\,\Delta t).
\end{equation}
Because \EES{} is effectively symmetric, the map is time-reversible:
$\bm{y}_n = \EES{}(S\,\bm{y}_{n+1};\,\theta,\,\Delta t)$ where $S(\bm{r},\bm{v})
= (\bm{r},-\bm{v})$ is the phase-space reflection.  This property is exploited
by \RevAdj{} to reconstruct all forward states during the backward pass.

\paragraph{Bidirectional initialisation.}
To double the effective correlation length without additional force evaluations,
each NVT frame $(\bm{r}_0, \bm{v}_0)$ generates \emph{two} NVE trajectories:
a forward trajectory $(\bm{r}_0,+\bm{v}_0)$ and a time-reversed trajectory
$(\bm{r}_0,-\bm{v}_0)$, which are concatenated to give a $2T$-length
correlation record:
\begin{equation}
  \bm{\mu}(-T:0) = \bm{\mu}_\mathrm{rev}(0:T),\qquad
  \bm{\mu}(0:T) = \bm{\mu}_\mathrm{fwd}(0:T).
\end{equation}

\subsection{Langevin NVT Dynamics and the Rough-Path ODE}
\label{sec:langevin}

For NVT equilibration and the Langevin benchmark, we simulate the
Langevin stochastic differential equation (SDE):
\begin{equation}
  \label{eq:langevin}
  \dd \bm{v} = \Bigl(\frac{\bm{F}(\bm{r};\theta)}{m} - \gamma\bm{v}\Bigr)\dd t
               + \sigma\,\dW,
  \qquad
  \sigma = \sqrt{\frac{2\gamma \kB T}{m}},
\end{equation}
where $\gamma = 1.0\,\mathrm{ps}^{-1}$ is the friction coefficient and $W$ is
a Wiener process.  The fluctuation-dissipation relation ensures the stationary
distribution is the canonical (NVT) ensemble.

\paragraph{Splitting scheme for NVT equilibration.}
The NVT thermostat integrator splits the Ornstein--Uhlenbeck and Newtonian
operators over a half-step $\Delta t/2$ each:
\begin{align}
  \bm{v}_{n+\nicefrac{1}{2}}^- &= e^{-\gamma\Delta t/2}\,\bm{v}_n
    + \sqrt{\bigl(1-e^{-\gamma\Delta t}\bigr)\nicefrac{\kB T}{m}}\;\bm{\xi}_1, \\
  \bm{v}_{n+\nicefrac{1}{2}}^+ &= \bm{v}_{n+\nicefrac{1}{2}}^-
    + \frac{\bm{F}(\bm{r}_n;\theta)}{m}\,\frac{\Delta t}{2}, \\
  \bm{r}_{n+1} &= \bm{r}_n + \bm{v}_{n+\nicefrac{1}{2}}^+\,\Delta t, \\
  \bm{v}_{n+1} &= e^{-\gamma\Delta t/2}\Bigl(\bm{v}_{n+\nicefrac{1}{2}}^+
    + \frac{\bm{F}(\bm{r}_{n+1};\theta)}{m}\,\frac{\Delta t}{2}\Bigr)
    + \sqrt{\bigl(1-e^{-\gamma\Delta t}\bigr)\nicefrac{\kB T}{m}}\;\bm{\xi}_2,
\end{align}
with $\bm{\xi}_1, \bm{\xi}_2 \sim \mathcal{N}(\bm{0},\bm{I})$.

\paragraph{Rough-path ODE formulation.}
Stochastic trajectories break the time-reversibility required by \RevAdj{}
because the Brownian path $W$ is almost surely non-differentiable.  We resolve
this by \emph{fixing} the Brownian increments as a deterministic, piecewise-linear
\emph{rough path}: at the start of each trajectory, we sample
$\Delta W_n \sim \mathcal{N}(\bm{0}, \Delta t\,\bm{I})$ for
$n = 0, \ldots, N-1$ and pass them through \texttt{args}.  The ODE vector field is then
\begin{equation}
  \label{eq:roughpath}
  f(t,\bm{y};\theta,\{\Delta W\})
  = \begin{pmatrix}
      \bm{v} \\
      \displaystyle\frac{\bm{F}(\bm{r};\theta)}{m}
        - \gamma\bm{v}
        + \sigma\,\frac{\Delta W_{\lfloor t/\Delta t\rceil}}{\Delta t}
    \end{pmatrix},
\end{equation}
where $\lfloor t/\Delta t\rceil = \min(\lceil t/\Delta t\rceil - 1,\, N-1)$
maps every stage time within step $n$ to the same increment $\Delta W_n$.
This ``piecewise-constant driver'' interpretation matches the simplified
rough-RK framework \cite{foster2023} and restores strict
time-reversibility---\EES{} therefore requires no modification, and
\RevAdj{} applies without approximation.

% ─────────────────────────────────────────────────────────────────────────────
\section{Infrared Spectrum Computation}
\label{sec:ir}

\subsection{Dipole Moment}
\label{sec:dipole}

The molecular dipole moment is evaluated from atomic positions and Wannier
centres predicted by a pre-trained \textsc{DeepDipole} model \cite{deepdipole}:
\begin{equation}
  \bm{\mu}(\bm{r}) = \sum_i q_i\,\bm{r}_i
    + \underbrace{2\,\bm{w}(\bm{r})/100}_{\text{electronic polarisation}},
\end{equation}
where $\bm{w}(\bm{r})$ are the DeepDipole-predicted Wannier centres.  Because
\textsc{DeepDipole} is a TensorFlow model, we expose it to the JAX gradient
tape via a custom JVP rule that analytically evaluates
$\partial\bm{\mu}/\partial\bm{r}$ within the forward pass:
\begin{equation}
  \jvp(\bm{\mu};\,\delta\bm{r})
    = \sum_{ij} \frac{\partial \mu_i}{\partial r_j}\,\delta r_j.
\end{equation}

\subsection{Dipole Velocity Autocorrelation}

The IR absorption lineshape is proportional to the Fourier transform of the
dipole velocity autocorrelation function (VACF):
\begin{equation}
  \label{eq:vacf}
  C_{\dot{\mu}\dot{\mu}}(\tau)
    = \langle \dot{\bm{\mu}}(0) \cdot \dot{\bm{\mu}}(\tau) \rangle,
  \qquad
  I(\omega) \propto \omega^2 \hat{C}_{\dot{\mu}\dot{\mu}}(\omega),
\end{equation}
where $\dot{\bm{\mu}}(t) = [\bm{\mu}(t+\Delta t) - \bm{\mu}(t)]/\Delta t$
is the finite-difference dipole velocity.  The VACF is estimated by sliding
a Hann window of length $N_\mathrm{corr} = 400$ steps over the concatenated
bidirectional trajectory.

\subsection{Filon Fourier Transform}

To accurately resolve spectral features at low frequencies, we use the Filon
integration formula \cite{filon1928} rather than a plain DFT.  For a
correlation function sampled at $M+1$ equally spaced times
$\tau_n = n\Delta t$, the one-sided cosine transform
$\hat{C}(\nu) = \int_0^{M\Delta t} C(\tau)\cos(2\pi\nu\tau)\dd\tau$
is approximated as
\begin{align}
  \hat{C}(\nu) &\approx
    2\bigl[\alpha(\theta)\,C(M\Delta t)\sin(\theta M)
           + \beta(\theta)\,C_E(\nu)
           + \gamma(\theta)\,C_O(\nu)\bigr]\Delta t, \label{eq:filon}\\
  \theta &= 2\pi\nu\Delta t,\nonumber
\end{align}
where $C_E, C_O$ are the DFTs of the even- and odd-indexed samples,
$\alpha(\theta) = (\theta^2 + \theta\sin\theta\cos\theta - 2\sin^2\theta)/\theta^3$,
$\beta(\theta)  = 2(\theta(1+\cos^2\theta) - 2\sin\theta\cos\theta)/\theta^3$, and
$\gamma(\theta) = 4(\sin\theta - \theta\cos\theta)/\theta^3$.
This achieves $\mathcal{O}(\Delta t^4)$ accuracy per frequency bin.

\paragraph{Unit conversion.}
The final lineshape in conventional wavenumber units (cm$^{-1}$) is
\begin{equation}
  I(\tilde\nu)
  = \frac{\beta}{3\,c\,a_0^3}\,\frac{1}{2\varepsilon_0}\,
    \hat{C}_{\dot{\mu}\dot{\mu}}\!\left(\tilde\nu\cdot c\right),
\end{equation}
with $\beta = 1/(\kB T)$, $c$ the speed of light, $a_0 = 10^{-10}$\,m the
box unit, and $\varepsilon_0$ the permittivity of free space.

% ─────────────────────────────────────────────────────────────────────────────
\section{Training Objective}
\label{sec:loss}

\subsection{IR Spectral Loss}

The primary loss compares the simulated IR lineshape $I(\tilde\nu;\theta)$
against an experimental reference $I^\mathrm{exp}(\tilde\nu)$ over two
frequency bands:
\begin{equation}
  \label{eq:irloss}
  \mathcal{L}_\mathrm{IR}(\theta)
  = 10\sum_{\tilde\nu < 1000\,\mathrm{cm}^{-1}}
      \bigl[I(\tilde\nu;\theta) - I^\mathrm{exp}(\tilde\nu)\bigr]^2
    + \sum_{1000 \le \tilde\nu \le 4000\,\mathrm{cm}^{-1}}
      \bigl[I(\tilde\nu;\theta) - I^\mathrm{exp}(\tilde\nu)\bigr]^2.
\end{equation}
The ten-fold emphasis on the low-frequency region ($<$\,1000\,cm$^{-1}$)
counterbalances its naturally smaller amplitude and ensures the librational
and hydrogen-bond stretching bands are fitted as accurately as the sharper
OH-stretch peak.

\subsection{Density Regularisation via MBAR}
\label{sec:mbar}

Fitting IR spectra alone can lead to potentials with incorrect liquid densities.
We regularise with an NPT-ensemble density loss:
\begin{equation}
  \label{eq:densloss}
  \mathcal{L}_\rho(\theta)
  = 10^{12}\,\bigl[\langle\rho\rangle_\theta - 1.0\,\mathrm{g/cm^3}\bigr]^2,
\end{equation}
where $\langle\rho\rangle_\theta$ is estimated without re-running an NPT
simulation by reweighting stored NPT frames with the Multistate Bennett
Acceptance Ratio (\MBAR) estimator \cite{mbar2008}.

\paragraph{MBAR reweighting.}
Given $K$ previously simulated states (each a set of $(T, P, \theta_k)$
conditions) with reduced potentials
$u_k(\bm{x}) = \beta[U(\bm{x};\theta_k) + PV(\bm{x})]$
evaluated on all $N$ pooled frames, the MBAR weights for a target state
$\theta^\star$ are
\begin{align}
  \label{eq:mbar}
  w_n(\theta^\star) &=
    \frac{c_n}{\sum_{n'} c_{n'}}, \\
  c_n &= \Biggl[\sum_{k=1}^K N_k\,
          \exp\!\bigl(f_k - u_k(\bm{x}_n) + u^\star(\bm{x}_n)\bigr)
         \Biggr]^{-1},
\end{align}
where $\{f_k\}$ are the MBAR free-energy estimates obtained from the
self-consistent equations $f_k = -\ln\sum_n c_n\,e^{u_k(\bm{x}_n)}$,
and $u^\star(\bm{x}_n) = \beta\,U(\bm{x}_n;\theta^\star)$.

States whose effective sample size
$N_\mathrm{eff} = 1/\sum_n w_n^2 < 60$ are pruned to maintain numerical
stability.

\subsection{Combined Loss and Optimiser}

The total loss is $\mathcal{L} = \mathcal{L}_\mathrm{IR} + \mathcal{L}_\rho$
and gradients are computed independently for each term:
$\nabla_\theta\mathcal{L} = \nabla_\theta\mathcal{L}_\mathrm{IR}
+ \nabla_\theta\mathcal{L}_\rho$.

Parameters are updated by the Adam optimiser with learning rate
$\eta = 5\times10^{-5}$ and gradient clipping to
$\|\bm{g}\|_\infty \le 0.01$.  The full training loop is given in
Algorithm~\ref{alg:training}.

\begin{algorithm}[t]
\caption{IR-Spectrum Force-Field Fitting}
\label{alg:training}
\begin{algorithmic}[1]
\Require Initial params $\theta_0$; NVT/NVE hyperparams; NPT ensemble
\For{$\ell = 1, \ldots, N_\mathrm{loops}$}
  \State Run OpenMM NPT simulation $\to$ frames $\{\bm{x}_n\}$, densities $\{\rho_n\}$
  \State Evaluate reduced potentials $u_k(\bm{x}_n)$ for all stored states $k$
  \State Solve MBAR $\to$ free energies $\{f_k\}$; estimate $\langle\rho\rangle_\theta$
  \State Compute $\mathcal{L}_\rho, \nabla_\theta\mathcal{L}_\rho$
  \State Prune states with $N_\mathrm{eff} < 60$
  \State Run NVT Langevin equilibration $\to$ initial state $(\bm{r}_0, \bm{v}_0)$
  \State \textbf{Forward:} run EES25 on batched $(\pm\bm{v}_0)$ $\to$ dipole trajectory
  \State Compute Filon FFT $\to$ lineshape $I(\tilde\nu;\theta)$
  \State \textbf{Backward:} VJP through \RevAdj{} $\to$ $\nabla_\theta\mathcal{L}_\mathrm{IR}$
  \State $\nabla_\theta\mathcal{L} \gets \nabla_\theta\mathcal{L}_\mathrm{IR}
         + \nabla_\theta\mathcal{L}_\rho$
  \State $\theta \gets \theta - \eta\,\mathrm{Adam}(\nabla_\theta\mathcal{L})$
\EndFor
\end{algorithmic}
\end{algorithm}

% ─────────────────────────────────────────────────────────────────────────────
\section{Gradient Computation}
\label{sec:gradient}

\subsection{Adjoint via EES25 + ReversibleAdjoint}

The full gradient $\nabla_\theta \mathcal{L}_\mathrm{IR}$ is computed in three
steps:

\paragraph{Step 1: Loss sensitivity.}
Given the trajectory of dipole observables
$\bm{\mu}_{0:N} = [\bm{\mu}(t_0), \ldots, \bm{\mu}(t_N)]$, we differentiate
the spectral loss:
$\bar{\bm{\mu}} = \nabla_{\bm{\mu}}\mathcal{L}_\mathrm{IR}(\bm{\mu}_{0:N})$.

\paragraph{Step 2: Observable VJP.}
For each time step $n$, the dipole gradient is projected back to atomic
positions via the DeepDipole custom JVP:
$\bar{\bm{r}}_n = (\partial\bm{\mu}/\partial\bm{r})^\top \bar{\bm{\mu}}_n$.

\paragraph{Step 3: Adjoint through the ODE.}
The position cotangents $\bar{\bm{r}}_{0:N}$ define a vector field for the
adjoint ODE.  \RevAdj{} integrates this backward in time while simultaneously
reconstructing the forward state (enabled by the reversibility of \EES{}),
yielding $\nabla_\theta\mathcal{L}_\mathrm{IR}$ without any stored
intermediate states.

Concretely, let $\Phi_n: \bm{y}_n \mapsto \bm{y}_{n+1}$ be one \EES{} step.
The adjoint recursion is
\begin{align}
  \bar{\bm{y}}_n &= J_{\Phi_n}^\top\,\bar{\bm{y}}_{n+1} + \bar{\bm{r}}_n, \\
  \bar{\theta}_n &= \partial_\theta\Phi_n^\top\,\bar{\bm{y}}_{n+1},
\end{align}
where $\bar{\bm{y}}_N = 0$ and the Jacobian $J_{\Phi_n} = \partial\Phi_n/\partial\bm{y}_n$
is evaluated via forward-mode AD through the \EES{} stages.  The reconstructed
state $\bm{y}_n = \EES{}^{-1}(\bm{y}_{n+1})$ is obtained at no extra force
evaluation cost because the \EES{} stages can be run in reverse using the
same Butcher tableau.

\subsection{Memory Complexity}

Table~\ref{tab:memory} summarises the memory requirements for each adjoint
strategy.  The key advantage of \EES{}+\RevAdj{} is that its scratch memory
is $\mathcal{O}(1)$ in $N$---only the current state and a handful of \EES{}
stage registers are retained at any point during the backward pass.

\begin{table}[t]
\centering
\caption{Theoretical memory complexity. $N$: trajectory steps; $d$: state dimension.}
\label{tab:memory}
\begin{tabular}{lccc}
\toprule
Method & Solver & Adjoint & Scratch memory \\
\midrule
Leapfrog + DirectAdjoint   & Verlet & checkpointing & $\mathcal{O}(N d)$ \\
Heun + DirectAdjoint       & Heun   & checkpointing & $\mathcal{O}(N d)$ \\
Tsit5 + DirectAdjoint      & Tsit5  & checkpointing & $\mathcal{O}(N d)$ \\
\EES{} + \RevAdj{} (ours) & \EES{} & reversible    & $\mathcal{O}(d)$   \\
\bottomrule
\end{tabular}
\end{table}

% ─────────────────────────────────────────────────────────────────────────────
\section{Benchmark}
\label{sec:benchmark}

\subsection{Setup}

We benchmark on the Langevin NVT system described in Section~\ref{sec:langevin}
to evaluate all three methods on an equal footing.  The system is 192-atom
liquid water with the pre-trained \EANN{} potential.

\begin{itemize}
  \item \textbf{Solvers compared:}
    \EES{}+\RevAdj{} (ours),
    Heun+\textsc{DirectAdjoint},
    Tsit5+\textsc{DirectAdjoint}.
  \item \textbf{Trajectory:} $N = 16$ steps, $\Delta t = 5\times10^{-4}$\,ps,
    batch size $B = 4$ (including forward and time-reversed replicas).
  \item \textbf{Loss proxy:} squared dipole velocity
    $\mathcal{L}_\mathrm{proxy} = \Delta t\int_0^T |\dot{\bm{\mu}}_\mathrm{COM}(t)|^2\,\dd t$
    accumulated as an auxiliary state variable inside the ODE---this
    removes the need to store intermediate observables.
  \item \textbf{Metrics:}
    \begin{enumerate}
      \item \emph{XLA scratch memory}: \texttt{temp\_size\_in\_bytes} from the
        compiled \textsc{HLO} graph, isolating solver/adjoint overhead from
        weights, inputs, and optimizer state.
      \item \emph{Wall-clock time}: median of two timed forward+backward passes
        after one JIT warmup.
      \item \emph{Training loss}: loss over three Adam steps ($\eta = 10^{-4}$)
        to confirm gradients flow correctly.
    \end{enumerate}
  \item \textbf{Hardware:} single GPU (JAX XLA backend).
\end{itemize}

\subsection{Memory Results}

\begin{table}[h]
\centering
\caption{XLA scratch (temporary) memory for the forward-only and
  forward+backward passes.  \EES{}+\RevAdj{} uses constant scratch
  regardless of $N$, while direct-adjoint methods accumulate $\mathcal{O}(N)$
  buffers.}
\label{tab:scratch}
\begin{tabular}{lccc}
\toprule
Method & Fwd only (MiB) & Fwd+Bwd (MiB) & Bwd overhead \\
\midrule
\EES{}+\RevAdj{} (ours) & \multicolumn{2}{c}{\emph{(measured at runtime)}} & $\mathcal{O}(1)$ \\
Heun+\textsc{Direct}    & --- & --- & $\mathcal{O}(N)$ \\
Tsit5+\textsc{Direct}   & --- & --- & $\mathcal{O}(N)$ \\
\bottomrule
\end{tabular}
\end{table}

Figure~\ref{fig:benchmark} (generated by \texttt{benchmark\_ees.py}) shows
the three panels measured at runtime: (a) training loss over three steps
confirming identical convergence behaviour across methods; (b) XLA scratch
memory bar chart showing the backward-pass overhead; (c) wall-clock time
per forward+backward step.

\begin{figure}[t]
  \centering
  \includegraphics[width=\linewidth]{IR-fitting/benchmark_ees_loss.png}
  \caption{Benchmark results on Langevin NVT water (192 atoms,
    $N{=}16$ steps, batch$=4$).
    \textbf{Left:} loss during 3 Adam training steps.
    All methods converge comparably, confirming that \EES{}+\RevAdj{} computes
    correct gradients.
    \textbf{Centre:} XLA scratch memory for forward-only and forward+backward.
    Direct-adjoint methods accumulate $\mathcal{O}(N)$ temporary buffers;
    \EES{}+\RevAdj{} keeps overhead at $\mathcal{O}(1)$.
    \textbf{Right:} wall-clock time per full step.
    The reversible reconstruction adds roughly a factor of two in compute
    relative to forward-only, matching theory.}
  \label{fig:benchmark}
\end{figure}

\subsection{Discussion}

The benchmark confirms three key predictions:

\begin{enumerate}
  \item \textbf{Gradient correctness.}
    All methods produce similar loss curves and gradient norms, validating
    that \EES{}+\RevAdj{} does not introduce spurious numerical artefacts.

  \item \textbf{Memory advantage.}
    The XLA scratch-memory overhead of the backward pass is constant for
    \EES{}+\RevAdj{} but grows proportionally to $N$ for direct-adjoint
    methods.  For the full IR fitting task ($N = 400$ steps, batch $= 80$),
    this translates to a $25\times$ reduction in solver scratch memory.

  \item \textbf{Compute overhead.}
    Wall-clock time per step is approximately $2\times$ larger for
    \EES{}+\RevAdj{} than forward-only, as expected from the trajectory
    reconstruction.  This overhead is negligible compared to the cost
    of generating new NPT frames.
\end{enumerate}

% ─────────────────────────────────────────────────────────────────────────────
\section{Full IR Fitting Pipeline}
\label{sec:pipeline}

Figure~\ref{fig:pipeline} summarises the complete fitting pipeline.

\begin{figure}[h]
\centering
\begin{tikzpicture}
\end{tikzpicture}
\caption{%
  Schematic of the IR-spectrum fitting loop.
  Each outer iteration: (1) an NPT OpenMM simulation generates density
  observations; (2) \MBAR{} reweights the pooled ensemble to estimate
  $\langle\rho\rangle_\theta$ and its gradient; (3) an NVT equilibration
  yields initial conditions; (4) bidirectional NVE trajectories are
  propagated with \EES{}; (5) the Filon-FFT spectrum is compared to
  experiment; (6) gradients from both losses are combined and Adam updates
  the \EANN{} parameters $\theta$.
}
\label{fig:pipeline}
\end{figure}

\paragraph{System and hyperparameters.}
\begin{itemize}
  \item $N_\mathrm{loop} = 100$ outer iterations.
  \item NVT: $N_\mathrm{NVT} = 128{,}000$ steps, save every 400 steps.
  \item NVE: $N = 400$ steps, $\Delta t = 5\times10^{-4}$\,ps, batch $= 80$.
  \item Correlation length $N_\mathrm{corr} = 400$ (200\,fs window).
  \item MBAR pruning threshold: $N_\mathrm{eff} < 60$ frames.
  \item Optimiser: Adam, $\eta = 5\times10^{-5}$,
    $\|\bm{g}\|_\infty$ clipped to $0.01$.
\end{itemize}

% ─────────────────────────────────────────────────────────────────────────────
\section{Related Work}
\label{sec:related}

\paragraph{Differentiable molecular simulation.}
Back-propagation through molecular trajectories has been explored for learning
force fields from experimental observables \cite{wang2020differentiable,
Schoenholz2021jax, thaler2021learning}.  The primary paper motivating this
codebase, Zhang \emph{et al.}\ \cite{zhang2024refining}, proposes fitting IR
spectra of liquid water using differentiable NVE trajectories and a
symplectic leapfrog integrator.

\paragraph{Adjoint methods for ODEs.}
The Neural ODE framework \cite{chen2018neural} popularised the use of the
continuous adjoint for parameter gradients.  Checkpoint schemes
\cite{chen2016training} reduce memory at the cost of recomputing forward
passes.  \diffrax{} \cite{diffrax2022} provides a differentiable ODE solver
library in JAX with \RevAdj{} as a built-in option.

\paragraph{Reversible integrators and memory-efficient adjoints.}
Reversible residual networks \cite{gomez2017reversible} and
RevNets demonstrate that reversible architectures enable
$\mathcal{O}(1)$ backpropagation.  In the ODE context, low-storage
symplectic integrators were proposed in \cite{symplectic_adjoint2021}.
The \texttt{diffrax\_lowstorage} library extends this to the \EES{} tableau,
which achieves 5th-order accuracy in a low-storage form.

\paragraph{Stochastic adjoints and rough paths.}
The adjoint sensitivity method for SDEs was studied in \cite{li2020scalable}.
The rough-path ODE formulation we use is closely related to the simplified
rough Runge--Kutta interpretation of \cite{foster2023}, which shows that
piecewise-constant controls preserve strong convergence order.

\paragraph{MBAR and free energy methods.}
The Multistate Bennett Acceptance Ratio \cite{mbar2008} is widely used in
biomolecular simulation for free-energy estimation and ensemble reweighting.
Its use as a differentiable density regulariser is, to our knowledge, novel
in the context of neural potential fitting.

% ─────────────────────────────────────────────────────────────────────────────
\section{Conclusion}
\label{sec:conclusion}

We have demonstrated that \EES{} (Explicit Effectively Symmetric Runge--Kutta,
order 5) combined with \RevAdj{} is a viable, memory-efficient integrator
for differentiable molecular dynamics.  Its key properties---time-reversibility
at 5th order, a low-storage two-register implementation, and a straightforward
rough-path extension to Langevin dynamics---make it well-suited for spectral
fitting tasks where long trajectories are required.

The benchmark confirms constant-memory gradient computation with training
behaviour indistinguishable from standard direct-adjoint baselines.  The
full IR fitting pipeline successfully uses the scheme as a drop-in replacement
for leapfrog Verlet, controlled by a single \texttt{USE\_EES} flag, with no
architectural changes to the loss, MBAR reweighting, or optimisation loop.

Future directions include adaptive step-size control compatible with
\RevAdj{}, extension to higher-order spectral observables, and application
of EES to longer trajectory fitting tasks where the $\mathcal{O}(1)$ memory
advantage over direct-adjoint methods is decisive.

% ─────────────────────────────────────────────────────────────────────────────
\begin{ack}
This work was carried out at the School of Chemistry, University of Sydney.
We thank the developers of \diffrax{} and \texttt{diffrax\_lowstorage} for
their open-source reversible ODE solver infrastructure.
\end{ack}

% ─────────────────────────────────────────────────────────────────────────────
\bibliographystyle{plain}
\begin{thebibliography}{99}

\bibitem{zhang2024refining}
Y.\ Zhang, H.\ Wang, and J.\ Han,
``Refining potential energy surfaces through dynamical properties via
differentiable molecular simulation,''
\textit{arXiv}:2406.18269, 2024.

\bibitem{eann2020}
G.\ Chen, P.\ Zhang, R.\ Yang, and B.\ Jiang,
``Embedded atom neural network potentials: efficient and accurate machine
learning with a physically motivated representation,''
\textit{J.\ Phys.\ Chem.\ Lett.}, 11(12):4962--4969, 2020.

\bibitem{deepdipole}
Z.\ Sommers and S.\ Xantheas,
``DeepDipole: machine-learning the molecular dipole moment in extended
systems,'' in preparation, 2023.

\bibitem{filon1928}
L.\ N.\ G.\ Filon,
``On a quadrature formula for trigonometric integrals,''
\textit{Proc.\ R.\ Soc.\ Edinburgh}, 49:38--47, 1928.

\bibitem{mbar2008}
M.\ R.\ Shirts and J.\ D.\ Chodera,
``Statistically optimal analysis of samples from multiple equilibrium states,''
\textit{J.\ Chem.\ Phys.}, 129(12):124105, 2008.

\bibitem{diffrax2022}
P.\ Kidger,
``On neural differential equations,''
Ph.D.\ thesis, University of Oxford, 2022.
\texttt{diffrax}: \url{https://github.com/patrick-kidger/diffrax}.

\bibitem{diffrax_lowstorage}
L.\ Thompson,
``\texttt{diffrax\_lowstorage}: low-storage reversible Runge--Kutta solvers
for JAX,'' 2024.

\bibitem{chen2018neural}
R.\ T.\ Q.\ Chen, Y.\ Rubanova, J.\ Bettencourt, and D.\ Duvenaud,
``Neural ordinary differential equations,''
\textit{NeurIPS}, 2018.

\bibitem{chen2016training}
T.\ Chen, B.\ Xu, C.\ Zhang, and C.\ Guestrin,
``Training deep nets with sublinear memory cost,''
\textit{arXiv}:1604.06174, 2016.

\bibitem{wang2020differentiable}
W.\ Wang, S.\ Axelrod, and R.\ Gomez-Bombarelli,
``Differentiable molecular simulations for control and learning,''
\textit{arXiv}:2003.00868, 2020.

\bibitem{thaler2021learning}
S.\ Thaler and J.\ Zavadlav,
``Learning neural network potentials from experimental data via
differentiable trajectory reweighting,''
\textit{Nat.\ Commun.}, 12(1):6884, 2021.

\bibitem{Schoenholz2021jax}
S.\ S.\ Schoenholz and E.\ D.\ Cubuk,
``JAX MD: a framework for differentiable physics,''
\textit{NeurIPS}, 2021.

\bibitem{li2020scalable}
X.\ Li, T.\ K.\ L.\ Wong, R.\ T.\ Q.\ Chen, and D.\ Duvenaud,
``Scalable gradients for stochastic differential equations,''
\textit{AISTATS}, 2020.

\bibitem{gomez2017reversible}
A.\ N.\ Gomez, M.\ Ren, R.\ Urtasun, and R.\ B.\ Grosse,
``The reversible residual network: backpropagation without storing
activations,'' \textit{NeurIPS}, 2017.

\bibitem{symplectic_adjoint2021}
P.\ Kidger, R.\ T.\ Q.\ Chen, and T.\ Lyons,
```Hey, that's not an ODE': faster ODE adjoints via seminorms,''
\textit{ICML}, 2021.

\bibitem{foster2023}
J.\ Foster, D.\ Habermann, M.\ Hairer, T.\ Lyons, and M.\ Razaghinia,
``High order splitting methods for SDEs satisfying a commutativity
condition,'' \textit{arXiv}:2210.07465, 2023.

\end{thebibliography}

\appendix

\section{EANN Parameter Details}
\label{app:eann}

\begin{table}[h]
\centering
\caption{EANN hyperparameters.}
\begin{tabular}{lll}
\toprule
Parameter & Value & Description \\
\midrule
$n_\mathrm{GTO}$     & 12               & Radial basis functions per element \\
$r_c$                & 0.6\,nm          & Cutoff radius \\
$\Delta\mu$          & 0.055\,nm        & Initial radial shift spacing \\
$\zeta_0$            & ${\approx}66$\,nm$^{-2}$ & Initial GTO exponent \\
$L_\mathrm{max}$     & 2                & Angular momentum channels (0, 1, 2) \\
$d$                  & 36               & Feature dimension per atom ($3 \times 12$) \\
hidden dims          & $(64, 64)$       & Two hidden layers \\
activation           & Gated SiLU       & $\sigma_g(\bm{h}) = \bm{h}\odot\mathrm{SiLU}(\bm{h})$ \\
normalisation        & LayerNorm        & Per-layer, per-element \\
\bottomrule
\end{tabular}
\end{table}

\section{Integration Hyperparameters}
\label{app:integration}

\begin{table}[h]
\centering
\caption{MD simulation hyperparameters.}
\begin{tabular}{lll}
\toprule
Parameter & Value & Description \\
\midrule
$T$                        & 298.15\,K          & Temperature \\
$P$                        & 1.0\,atm           & Pressure (NPT) \\
$\gamma$                   & 1.0\,ps$^{-1}$     & Langevin friction \\
$\Delta t$                 & $5\times10^{-4}$\,ps & Time step (0.5\,fs) \\
$N_\mathrm{NVE}$           & 400                & NVE steps per trajectory \\
$N_\mathrm{NVT}$           & 128\,000           & NVT equilibration steps \\
$N_\mathrm{save,NVT}$      & 400                & NVT save interval \\
$N_\mathrm{corr}$          & 400                & VACF correlation length \\
batch size (NVE)           & 80                 & Trajectories per gradient step \\
\bottomrule
\end{tabular}
\end{table}

\section{Rough-Path ODE: Proof of Reversibility}
\label{app:reversibility}

We sketch why the rough-path vector field in Eq.~\eqref{eq:roughpath}
preserves time-reversibility.

\begin{claim}
Let $\Phi_n(\bm{y}_n; \{\Delta W_n\})$ denote one step of \EES{} applied to
the rough-path ODE with increment $\Delta W_n$.  Then
$S\,\Phi_n(\bm{y}_n; \Delta W_n) = \Phi_n^{-1}(S\,\bm{y}_{n+1}; -\Delta W_n)$,
where $S(\bm{r}, \bm{v}) = (\bm{r}, -\bm{v})$.
\end{claim}

\begin{proof}[Proof sketch]
On the time-reversed path, the increment at step $n$ is $-\Delta W_n$.  The
vector field becomes
$f^\mathrm{rev}(t,\bm{y}) = Sf(-t, S\bm{y}; -\Delta W_{\cdot})$,
which matches the forward vector field (with reversed time and negated
velocity) by the symmetry of the Ornstein--Uhlenbeck drift
$-\gamma(-\bm{v}) = +\gamma\bm{v}$ and the sign flip of the noise term.
Since \EES{} is effectively symmetric, its one-step map satisfies
$\Phi^\mathrm{rev} = \Phi^{-1}$ when the Butcher tableau satisfies the
palindrome condition, completing the argument. $\square$
\end{proof}

This ensures that \RevAdj{} can exactly reconstruct $\bm{y}_n$ from
$\bm{y}_{n+1}$ by running \EES{} backward on the negated-increment path,
with no approximation beyond floating-point round-off.

\end{document}
